\centerline{\bf \Huge Всероссная тренировка III}

\begin{flushright}
        \scriptsize{}
\end{flushright}

\problem{1}
Найдите количество корней уравнения

$$
|x|+|x+1|+\ldots+|x+2018|=x^2+2018 x-2019
$$


\problem{2}
Дано нечетное натуральное число $n>1$. На доске записаны числа $n, n+1, n+2, \ldots, 2 n-1$. Докажите, что можно стереть одно из них так, чтобы сумма оставшихся чисел не делилась ни на одно из оставшихся чисел.

\problem{3}
На столе по кругу разложены 1000 карточек, на каждой написано по натуральному числу; все эти числа различны. Сначала Вася выбирает одну из карточек и снимает её со стола. Далее он повторяет следующую операцию. Если на последней снятой карточке было написано число $k$, то Вася отсчитывает от неё по часовой стрелке $k$-ую не снятую со стола карточку и тоже снимает её. Это происходит до тех пор, пока на столе не останется одна карточка. Могло ли оказаться, что в начальном расположении есть такая карточка $A$, что если снять первой любую другую карточку, то в конце останется обязательно карточка $A$?

\problem{4}
Дан остроугольный треугольник $A B C$, в котором $A B<A C$. Пусть $M$ и $N$ - середины сторон $A B$ и $A C$ соответственно, а $D$ --- основание высоты, проведённой из $A$. На отрезке $M N$ нашлась точка $K$ такая, что $B K=C K$. Луч $K D$ пересекает окружность $\Omega$, описанную около треугольника $A B C$, в точке $Q$. Докажите, что точки $C, N, K$ и $Q$ лежат на одной окружности.